\section{Mẫu CI/CD hoàn chỉnh}
\label{sec:pc-tongquan}

Phụ lục này trình bày hai mẫu tích hợp liên tục và chuyển giao liên tục hoàn chỉnh cho dự án ASP.NET Web Forms trên .NET Framework 4.8.1.
Hai mẫu này đã được rút gọn cho khuôn khổ bài trình bày của nhóm nhưng vẫn giữ đầy đủ các bước quan trọng nhất trong thực tế.
Bạn đọc có thể sao chép trực tiếp hai tập tin bên dưới vào kho mã nguồn của mình và điều chỉnh tên solution cho phù hợp.
Toàn bộ ví dụ đều giả định solution có tên \verb|AspxDemo_NET481.sln| và dự án kiểm thử có tên \verb|AspxDemo.Tests|.
Các mẫu dưới đây tuân thủ tài liệu chính thức của GitHub Actions \cite{ghactions,githubdocs} và GitLab CI \cite{gitlabci,gitlabdocs}.
Quy trình build dựa trên MSBuild và bộ công cụ Visual Studio \cite{msmsbuild,msnetfx}.

Luồng chung của cả hai mẫu gồm ba giai đoạn nối tiếp nhau.
Giai đoạn thứ nhất biên dịch solution ở cấu hình Release trên máy Windows.
Giai đoạn thứ hai chạy bộ kiểm thử đơn vị bằng công cụ \verb|vstest.console.exe|.
Giai đoạn thứ ba đóng gói thư mục publish thành artifact để tải về hoặc triển khai tiếp.
Hình \ref{fig:pc-luong} minh họa luồng ba giai đoạn này một cách khái quát.
Việc tách thành ba giai đoạn giúp xác định nhanh bước nào bị lỗi khi pipeline báo đỏ.

\begin{figure}[H]
\centering
\begin{tikzpicture}[>=Stealth, thick, box/.style={draw, rounded corners, fill=blue!10, minimum width=2.4cm, minimum height=0.9cm, align=center, font=\small}]
\node[box] (a) {Build};
\node[box, right=1.2cm of a] (b) {Test};
\node[box, right=1.2cm of b] (c) {Package};
\draw[->] (a) -- node[midway, above, font=\footnotesize] {bien dich} (b);
\draw[->] (b) -- node[midway, above, font=\footnotesize] {kiem thu} (c);
\end{tikzpicture}
\caption{Luong CI/CD ba giai doan cho du an ASPX}
\label{fig:pc-luong}
\end{figure}

Khi áp dụng vào đồ án môn học, nhóm nên bắt đầu từ mẫu GitHub Actions vì chỉ cần dùng máy chủ có sẵn của GitHub.
Sau khi đã quen, nhóm có thể sao chép cùng ý tưởng sang GitLab CI để so sánh sự khác nhau về cú pháp.
Cả hai mẫu đều dùng máy Windows vì .NET Framework 4.8.1 không biên dịch được trên máy Linux thông thường.
Lệnh \verb|dotnet build| dành cho .NET hiện đại không thay thế được MSBuild trong trường hợp này.

\subsection{GitHub Actions}
\label{subsec:pc-github}

Mẫu trong mục này đặt tại đường dẫn \verb|.github/workflows/ci-aspx-481.yml| trong kho trên GitHub.
Quy trình được kích hoạt mỗi khi có lệnh đẩy lên nhánh \verb|main| hoặc \verb|develop| và mỗi khi có pull request nhắm vào hai nhánh này.
Máy chạy được chọn là \verb|windows-latest| để có sẵn Visual Studio, MSBuild, NuGet và công cụ kiểm thử.
Danh sách ví dụ Lst. \ref{lst:pc-github} trình bày đầy đủ nội dung tập tin workflow.
Các chú thích trong khối mã đều viết bằng tiếng Anh không dấu để giữ đúng quy tắc mã nguồn ASCII.

\begin{lstlisting}[language=yaml, caption={Mau GitHub Actions cho ASPX 4.8.1}, label={lst:pc-github}]
# CI for ASP.NET Web Forms 4.8.1 on GitHub Actions
# File: .github/workflows/ci-aspx-481.yml
name: ci-aspx-481
on:
  push:
    branches: [main, develop]
  pull_request:
    branches: [main, develop]
env:
  SOLUTION: AspxDemo_NET481.sln
  CONFIG: Release
  PLATFORM: "Any CPU"
jobs:
  build:
    runs-on: windows-latest
    steps:
      - name: Checkout code
        uses: actions/checkout@v4
      - name: Setup MSBuild
        uses: microsoft/setup-msbuild@v2
      - name: Setup NuGet
        uses: NuGet/setup-nuget@v2
      - name: Cache NuGet packages
        uses: actions/cache@v4
        with:
          path: ~/.nuget/packages
          key: "nuget-windows-${{ hashFiles('**/packages.config') }}"
          restore-keys: nuget-windows-
      - name: Restore packages
        run: nuget restore $env:SOLUTION
      - name: Build solution
        run: msbuild $env:SOLUTION /p:Configuration=$env:CONFIG /p:Platform="Any CPU" /m /v:minimal /t:Rebuild
      - name: Run unit tests
        run: vstest.console.exe .\Tests\bin\Release\AspxDemo.Tests.dll /Logger:trx
      - name: Publish web project
        run: msbuild .\Web\AspxDemo.csproj /p:Configuration=$env:CONFIG /p:DeployOnBuild=true /p:WebPublishMethod=FileSystem /p:PublishUrl=C:\Publish\ /v:minimal
      - name: Upload artifact
        uses: actions/upload-artifact@v4
        with:
          name: aspx-package
          path: C:\Publish\
\end{lstlisting}

Các bước trong workflow trên có thể hiểu như sau.

\begin{enumerate}
\item Bước \verb|Checkout code| tải toàn bộ mã nguồn của lần đẩy hiện tại về máy chạy.
\item Hai bước \verb|Setup MSBuild| và \verb|Setup NuGet| cài đặt công cụ build chuẩn của Microsoft vào đường dẫn hệ thống.
\item Bước \verb|Cache NuGet packages| lưu thư mục gói NuGet giữa các lần chạy để giảm thời gian tải gói.
Khóa cache được tính từ nội dung các tập tin \verb|packages.config| nên cache tự động mất hiệu lực khi đổi gói.
\item Bước \verb|Restore packages| khôi phục các thư viện bên thứ ba trước khi biên dịch.
Nếu thiếu bước này thì MSBuild sẽ báo lỗi không tìm thấy assembly tham chiếu.
\item Bước \verb|Build solution| biên dịch toàn bộ solution ở cấu hình Release với nhiều tiến trình song song nhờ cờ \verb|/m|.
Cờ \verb|/t:Rebuild| bảo đảm thư mục trung gian cũ bị xóa trước khi biên dịch lại từ đầu.
\item Bước \verb|Run unit tests| chạy assembly kiểm thử và xuất kết quả định dạng \verb|trx| để xem trên trang Actions.
Nếu bất kỳ kiểm thử nào thất bại thì toàn bộ workflow dừng lại và báo đỏ.
\item Hai bước cuối cùng xuất bản dự án web ra thư mục \verb|C:\Publish\| rồi tải thư mục này lên mục artifact.
Thành viên trong nhóm có thể tải artifact về để triển khai thủ công lên Internet Information Services.
\end{enumerate}

Khi mới bắt đầu, nhóm nên giữ nguyên tên nhánh \verb|main| và \verb|develop| như trong mẫu.
Sau này khi đã quen, nhóm có thể bổ sung thêm quy tắc bảo vệ nhánh yêu cầu workflow phải chạy xanh mới được hợp nhất.
Chuỗi kết nối cơ sở dữ liệu và các khóa bí mật tuyệt đối không ghi trực tiếp trong tập tin YAML.
Thay vào đó, nhóm khai báo secret trong mục Settings của kho rồi tham chiếu bằng biểu thức \verb|secrets.TEN_BI_MAT|.
Cách làm này giúp mã nguồn công khai nhưng thông tin nhạy cảm vẫn được bảo vệ an toàn.

\subsection{GitLab CI}
\label{subsec:pc-gitlab}

Mẫu trong mục này đặt tại tập tin \verb|.gitlab-ci.yml| ở thư mục gốc của dự án trên GitLab.
Pipeline gồm đúng ba stage là \verb|build|, \verb|test| và \verb|package| chạy nối tiếp nhau theo thứ tự khai báo.
Mỗi công việc đều gắn nhãn \verb|windows| trong mục \verb|tags| để GitLab điều phối đúng đến máy chạy Windows.
Yêu cầu này là bắt buộc vì máy chạy Linux không có MSBuild đầy đủ cho .NET Framework 4.8.1.
Danh sách ví dụ Lst. \ref{lst:pc-gitlab} trình bày đầy đủ nội dung tập tin cấu hình.

\begin{lstlisting}[language=yaml, caption={Mau GitLab CI cho ASPX 4.8.1}, label={lst:pc-gitlab}]
# CI for ASP.NET Web Forms 4.8.1 on GitLab CI
# File: .gitlab-ci.yml in repository root
stages:
  - build
  - test
  - package
variables:
  SOLUTION: AspxDemo_NET481.sln
  CONFIG: Release
  NUGET_PACKAGES: packages
cache:
  key: nuget-$CI_COMMIT_REF_SLUG
  paths:
    - packages/
before_script:
  - echo Building $env:SOLUTION
build:
  stage: build
  tags:
    - windows
  script:
    - nuget restore $env:SOLUTION
    - msbuild $env:SOLUTION /p:Configuration=$env:CONFIG /p:Platform="Any CPU" /m /v:minimal /t:Rebuild
  artifacts:
    paths:
      - Publish/
      - Tests/bin/Release/
    expire_in: 1 week
test:
  stage: test
  tags:
    - windows
  script:
    - vstest.console.exe .\Tests\bin\Release\AspxDemo.Tests.dll /Logger:"trx;LogFileName=results.trx"
    - dotnet tool install -g trx2junit; if ($LASTEXITCODE -ne 0) { echo "trx2junit already installed" }
    - trx2junit TestResults\results.trx --output TestResults\junit.xml
  artifacts:
    when: always
    reports:
      junit: TestResults/junit.xml
    paths:
      - TestResults/
package:
  stage: package
  tags:
    - windows
  script:
    - msbuild .\Web\AspxDemo.csproj /p:Configuration=$env:CONFIG /p:DeployOnBuild=true /p:WebPublishMethod=FileSystem /p:PublishUrl=Publish\ /v:minimal
  artifacts:
    paths:
      - Publish/
    expire_in: 1 month
\end{lstlisting}

Ý nghĩa của từng khối trong tập tin trên được giải thích chi tiết như sau.

\begin{itemize}
\item Khối \verb|stages| khai báo ba giai đoạn theo đúng thứ tự thực hiện trong pipeline.
Nếu một công việc ở stage trước thất bại thì các stage sau sẽ không được chạy.
\item Khối \verb|variables| định nghĩa các biến dùng chung như tên solution và cấu hình build.
Cách đặt biến tập trung giúp khi đổi tên solution chỉ cần sửa một chỗ duy nhất.
\item Khối \verb|cache| giữ lại thư mục gói NuGet giữa các lần chạy để pipeline chạy nhanh hơn.
Khóa cache dùng kèm tên nhánh nên mỗi nhánh có không gian cache riêng biệt.
\item Khối \verb|before_script| in ra tên solution đang build để dễ dàng kiểm tra nhật ký khi có nhiều dự án.
\item Công việc \verb|build| chạy trên máy Windows, khôi phục gói rồi biên dịch toàn bộ solution.
Thư mục \verb|Publish/| được giữ lại thành artifact trong một tuần để các stage sau tái sử dụng.
\item Công việc \verb|test| chạy assembly kiểm thử và luôn xuất báo cáo kể cả khi có kiểm thử thất bại nhờ \verb|when: always|.
Báo cáo định dạng \verb|trx| được GitLab hiển thị trực tiếp trong trang chi tiết của pipeline.
\item Công việc \verb|package| xuất bản dự án web ra thư mục \verb|C:\Publish\| và lưu lại trong một tháng.
Nhóm có thể tải gói này về để triển khai lên máy chủ Internet Information Services của khoa.
\end{itemize}

Để pipeline trên chạy được, dự án GitLab cần có ít nhất một máy chạy Windows đang hoạt động.
Máy chạy này có thể là máy tính của thành viên trong nhóm đăng ký ở chế độ shell hoặc máy chủ dùng chung của khoa.
Khi đăng ký máy chạy, nhóm phải gắn nhãn \verb|windows| cho đúng với mục \verb|tags| trong tập tin cấu hình.
Nếu nhãn không khớp thì pipeline sẽ mãi ở trạng thái chờ mà không có máy nào nhận việc.
Tài liệu GitLab CI mô tả chi tiết về stage, job, artifact và cache \cite{gitlabci,gitlabdocs}.

\subsection{Tham số MSBuild hay dùng}
\label{subsec:pc-msbuild}

Cả hai mẫu pipeline ở trên đều gọi MSBuild với một số tham số dòng lệnh giống nhau.
Bảng \ref{tab:pc-msbuild} tổng hợp tám tham số hay dùng nhất khi biên dịch solution ASP.NET Web Forms.
Nhóm nên thuộc ý nghĩa của tám tham số này trước buổi demo vì giảng viên thường hỏi về cờ build.
Tài liệu MSBuild liệt kê đầy đủ hàng trăm tham số khác cho các tình huống nâng cao \cite{msmsbuild}.

\begin{table}[H]
\centering
\caption{8 tham so MSBuild hay dung cho ASPX}
\label{tab:pc-msbuild}
\begin{tabular}{lll}
\toprule
Tham số & Ý nghĩa & Ví dụ \\
\midrule
\texttt{/p:Configuration} & Chọn cấu hình biên dịch & \texttt{/p:Configuration=Release} \\
\texttt{/p:Platform} & Chọn nền tảng đích & \texttt{/p:Platform="Any CPU"} \\
\texttt{/m} & Build song song nhiều tiến trình & \texttt{/m} \\
\texttt{/v} & Mức chi tiết nhật ký & \texttt{/v:minimal} \\
\texttt{/t:Rebuild} & Xóa rồi biên dịch lại từ đầu & \texttt{/t:Rebuild} \\
\texttt{/t:Clean} & Chỉ xóa thư mục trung gian & \texttt{/t:Clean} \\
\texttt{/p:DeployOnBuild} & Xuất bản ngay khi build & \texttt{/p:DeployOnBuild=true} \\
\texttt{/restore} & Khôi phục gói NuGet trước khi build & \texttt{/restore} \\
\bottomrule
\end{tabular}
\end{table}

Trong tám tham số trên, ba tham số \verb|/p:Configuration|, \verb|/p:Platform| và \verb|/m| xuất hiện trong hầu hết mọi pipeline.
Tham số \verb|/v:minimal| giúp nhật ký vừa đủ chi tiết mà không quá dài gây khó đọc trên trang web.
Cặp \verb|/t:Clean| và \verb|/t:Rebuild| đặc biệt hữu ích khi thư mục \verb|bin| hoặc \verb|obj| bị kẹt tập tin cũ.
Tham số \verb|/p:DeployOnBuild=true| thường đi kèm \verb|/p:WebPublishMethod| và \verb|/p:PublishUrl| để xuất bản web.
Tham số \verb|/restore| thay thế cho lệnh \verb|nuget restore| riêng lẻ trong các dự án mới dùng PackageReference.
Với dự án ASPX cũ dùng \verb|packages.config|, nhóm vẫn nên giữ bước \verb|nuget restore| riêng như trong hai mẫu trên để tránh lỗi thiếu assembly.

Tóm lại, phụ lục này cung cấp hai mẫu cấu hình sẵn sàng sao chép cho cả GitHub Actions và GitLab CI.
Nhóm chỉ cần đổi tên solution, tên dự án web và đường dẫn assembly kiểm thử là có pipeline chạy xanh.
Khi đã có pipeline xanh, nhóm nên bổ sung thêm kiểm tra chất lượng mã và quét bảo mật ở các bước tiếp theo.
Việc nắm vững tám tham số MSBuild trong Bảng \ref{tab:pc-msbuild} giúp nhóm tự tin giải thích trước giảng viên.
Nhóm nên lưu hai mẫu này vào kho tài liệu chung để tái sử dụng cho các đồ án tiếp theo.
